\section{Related Work}
\label{sec:related}

\para{LLMs as zero-shot rankers.} \citet{hou2024llm} formalize recommendation as conditional ranking and show that GPT-3.5, given a sequential interaction history and a $20$-track candidate pool, can rank competitively with classical sequential models on movie and book benchmarks; they also document position bias and popularity bias. \citet{dai2023uncovering} reach a similar conclusion for ChatGPT across multiple recommendation tasks. \citet{wang2023zero} study zero-shot next-item recommendation. We adopt Hou et al.'s sequential ranking prompt verbatim, but apply it for the first time to music histories under the multi-criterion evaluation framework of \citet{epure2025music}, with frontier $2025$-era LLMs.

\para{Cold-start and language-only recommendation.} \citet{sanner2023llms} compare LLMs against item-based collaborative filtering (EASE \citep{steck2019embarrassingly}) and BM25 on Amazon movie reviews and find LLMs match item-CF in the near cold-start regime even with only language preferences. Unlike their setting, our users are not cold-start: they have hundreds-to-thousands of historical interactions, which is the regime that platform recommenders actually serve.

\para{LLM-generated taste profiles.} \citet{radlinski2022natural} introduce natural-language user profiles for transparent recommendation. \citet{sguerra2025biases} systematically evaluate LLM-generated music taste profiles on $64$ Deezer users with three LLMs and show that self-identification and downstream recall are only weakly correlated; they also document item-side biases (rap content lowers profile quality, metal content raises it). Their finding motivates our choice to evaluate \emph{accuracy} on held-out tracks rather than subjective profile satisfaction.

\para{LLMs as end-to-end music recommenders.} The closest published precedent to our work is \citet{boadana2025llm}, who collect $13$ months of Spotify-export data from $19$ users and compare a CrewAI multi-agent system on Gemini~$2.0$ Flash and LLaMA-$3.3$-$70$B against a TF-IDF cosine content-based filter, with blind subjective ratings. Their LLaMA-based system achieves an $89\%$ like rate against $61\%$ for the CBF baseline, but trades off discovery (novelty rate $12\%$ \versus $58\%$) and pays a $50$--$60\times$ latency cost. Our work differs in three ways: we measure objective held-out-track accuracy (not subjective like rate), we use a closed-pool ranking protocol so that all methods rank identical candidates, and we compare against a broader baseline set including item-kNN, popularity, and random.

Generative retrieval approaches such as Text$2$Tracks \citep{palumbo2025text2tracks} and tool-using agents such as TalkPlay-Tools \citep{doh2025talkplay} represent more elaborate LLM-recommender architectures than the prompt-only setup we study; we treat them as complements rather than competitors, since our question is what off-the-shelf prompting (the user-accessible operational mode) achieves.

\para{Evaluation frameworks for generative recommenders.} \citet{epure2025music} synthesize a six-axis evaluation framework (G1-G6) covering query adherence, discovery, personalization gain, profile fidelity, cultural coverage, and classical relevance, plus risk diagnostics for popularity bias \citep{abdollahpouri2019popularity} and hallucination. We adopt G1, G2, G3, and G6 plus the popularity diagnostic. \citet{deldjoo2024recommendation} survey generative recommenders broadly. \citet{lin2023how} survey LLM-recommender integrations. None of these surveys provide the specific Claude-vs-classical comparison the user's question demands.

\para{Music recommendation datasets.} The Million Song Dataset \citep{bertin2011million}, $30$Music \citep{turrin201530music}, LFM-1b \citep{schedl2017lfm1b}, and Spotify Million Playlist (now no longer public) are the standard music datasets. We use the $50$-user sample of Last.fm 1K because (a) it has per-user time-stamped listens that mirror Spotify export semantics, and (b) it is the dataset most consistent with prior work in the LLM-recommender music subfield.

\para{Position relative to prior work.} Building on \citet{hou2024llm}, we apply zero-shot ranking to music; building on \citet{boadana2025llm}, we replace subjective like-rate with objective held-out accuracy and broaden the baseline set; building on \citet{epure2025music}, we report the four core diagnostic axes simultaneously on a single benchmark. Unlike all of these, we test two frontier LLMs (\claude and \gpt) to show that the result is not an artifact of any single model.
